% !TEX root = ../main.tex

\chapter{Introduction}
\label{ch:introduction}

\section{Problem description}
\label{sec:problem}
The evaluation of non-linear elementary functions, in particular trigonometric
functions such as the sine function, $\sin(x)$, is a frequent operation in
digital systems. Applications such as machine learning and digital signal
processing can often make do with fast, approximate evaluations of these
functions~\cite{muller2020}.
This project considers real-valued inputs represented in the 32-bit single-precision
(binary32) floating-point format defined by the IEEE~754 standard~\cite{ieee754}.

\section{Motivation}
\label{sec:motivation}
Elementary functions such as $\sin(x)$ cannot be computed exactly with a finite
number of arithmetic operations. Any implementation, either software or hardware,
can only approximate them~\cite{muller2020}. The precision of the result therefore
becomes a design decision. General-purpose implementations aim for results as
close as possible to the exact value, while many applications can tolerate a less
accurate result if it is obtained faster~\cite{muller2020}.

Choosing how to make this trade-off is not straightforward. In his review of
approximation techniques, Muller~\cite{muller2020} does not recommend a single
method: the most suitable one depends on the function being evaluated, the
required accuracy, the implementation technology and the memory available for
tables. He also notes that the methods differ in versatility, accuracy and
scalability, which makes a general comparison difficult.

On a Field-Programmable Gate Array (FPGA), this trade-off can be measured
directly. Each approximation method has a concrete cost in look-up tables (LUTs),
DSP slices, block RAM, latency and maximum clock frequency, and the best choice
can change with the required precision. For example, Muller reports work in which
CORDIC proved to be a good candidate for low-precision FPGA implementations of the
arctangent function~\cite{muller2020}.

General guidance from the literature is therefore not enough to select a method
for a specific platform and number format. A reliable choice requires the
candidate methods to be implemented on the same device and measured under
identical conditions, so that their accuracy and hardware cost can be compared
directly, including against the solutions already offered by the FPGA vendor.

\section{Objectives}
\label{sec:objectives}

The main objective of this project is to design, implement and evaluate two
pipelined hardware architectures that compute the sine of IEEE~754
single-precision floating-point numbers on an FPGA, exposed through AXI4-Stream
interfaces, and to compare their accuracy, resource usage and latency with the
Xilinx IP-cores available in the Vivado design suite.

The methods will be described in synthesizable VHDL, so that they can be
implemented on a real device and validated in hardware on a Digilent Nexys A7
development board, in addition to being verified in simulation.

\paragraph{Secondary objectives:}
\label{obj:secondary}
\begin{enumerate}[label=\textbf{\arabic*.}, ref=\arabic*, leftmargin=*]
  \item \label{obj:study}
    Study the existing methods for approximating elementary functions and
    select the two methods to be implemented: CORDIC and a table-based
    method.
  \item \label{obj:range}
    Design a range reduction stage that maps any finite binary32 input to
    the interval supported by the approximation methods, and handles the
    special values ($\pm 0$, subnormals, $\pm\infty$, NaN).
  \item \label{obj:cordic}
    Implement a pipelined CORDIC architecture for the sine function.
  \item \label{obj:lut}
    Implement a pipelined table-based architecture for the sine function.
  \item \label{obj:axi}
    Wrap both architectures in AXI4-Stream interfaces that respect the
    handshake rules.
  \item \label{obj:verify}
    Verify both designs in simulation against the software model, and
    measure their absolute and relative error over the binary32 input
    range.
  \item \label{obj:compare}
    Synthesize the designs together with the equivalent Vivado IP-cores
    and compare accuracy, resource usage (LUTs, flip-flops, DSP slices,
    block RAM), latency and maximum clock frequency.
  \item \label{obj:board}
    Demonstrate the designs running on the Nexys A7 board.
\end{enumerate}